\section{Zusammenfassung und Ausblick}
\label{sec:fazit}

\subsection{Zusammenfassung}
Die vorliegende Arbeit hat Potenziale und Herausforderungen Digitaler Zwillinge in der Fabrikplanung am Beispiel der Plattformen NVIDIA Omniverse und Isaac Sim systematisch untersucht. Auf Basis einer normativ fundierten Recherche, mit den Leitstandards ISO~23247, DIN~SPEC~91345 (RAMI~4.0), VDI~4499, VDI~5200, VDI/VDE~3693 und DIN~EN~IEC~62541, wurden für zwei reale Anlagen des LFP der HsH Digitale Zwillinge entwickelt und in Betrieb genommen. Der Rundtakttisch \enquote{Maze Runner} wurde via OPC~UA und MQTT, der Desktoproboter \enquote{Dobot} via ROS\,2 an Isaac Sim gekoppelt. Der Steuerpfad des Maze Runners erreicht dabei zuverlässig Latenzen zwischen \SI{100}{\milli\second} und \SI{500}{\milli\second}, während sich die Rückmeldung der Sensorwerte über MQTT im Laborbetrieb als nicht durchgängig zuverlässig erwies, sodass die Kopplung eher einer erweiterten Fernsteuerung als der vollständigen Klassifikation eines bidirektionalen Digitalen Zwillings nach Kritzinger et al.\ entspricht~\cite{kritzinger2018}. Der Dobot Magician erreicht mit dem automatisierten Datenfluss vom realen Roboter zur Simulation die Stufe des Digitalen Schattens und ergänzt diesen um eine manuell auslösbare Fernsteuerung in Gegenrichtung.
Die methodische Reproduzierbarkeit wurde durch einen siebenstufigen Erstellungsprozess sowie durch komplementäre Medienformate (PowerPoint, VR) sichergestellt. Aus den Ergebnissen wurden vier zentrale Nutzungspotenziale, virtuelle Inbetriebnahme, Kollisionsüberwachung, datengetriebene Erweiterbarkeit und KI-gestützte synthetische Datengenerierung, abgeleitet und konkrete Handlungsempfehlungen für das LFP formuliert.

\subsection{Ausblick}
Mehrere Anschlussarbeiten sind aus den Ergebnissen ableitbar:
\begin{itemize}
    \item \textbf{Skalierung auf weitere Anlagen:} Übertragung des Vorgehensmodells auf komplexere Produktionslinien und vernetzte Fertigungsverbünde.
    \item \textbf{Erprobung weiterer Manipulatoren:} Das von NVIDIA bereitgestellte Tutorial zu Import und Assemblierung von Manipulatoren in Isaac Sim ist auf die im LFP vorhandenen UR-Roboterarme und vergleichbare Greifersysteme übertragbar und bietet interessierten Personen einen vertiefenden Einstieg, ohne dass zunächst ein Hardwarebezug oder ein Digitaler Zwilling hergestellt werden muss~\cite{nvidia_isaacsim_manipulator_tutorial}.
    \item \textbf{Integration der Asset Administration Shell:} Anbindung der entwickelten Digitalen Zwillinge an eine AAS-konforme Datenstruktur gemäß IDTA \cite{idta_aas2025,zhang2025aas}.
    \item \textbf{KI-gestützte Modellgenerierung:} Einsatz von Large Language Models zur (teil-)automatisierten Erzeugung Digitaler Zwillinge aus textuellen Anlagen\-beschrei\-bungen, wie von Liu et al.\ \cite{mdpi2025isaac} für robotische Anwendungen demonstriert.
    \item \textbf{Industrial Metaverse:} Verknüpfung mehrerer Digitaler Zwillinge in einer gemeinsamen, kollaborativen Omniverse-Session zur standortübergreifenden Fabrikplanung.
    \item \textbf{Wirtschaftlichkeitsanalyse:} Quantitative Untersuchung der Kosten-Nutzen-Relation Digitaler Zwillinge in mittelständischen Produktionsunternehmen.
\end{itemize}
